We apply Huntington-Hill equal proportions apportionment to Census Bureau medium-series
2030 state population projections to produce the first systematic forecast of 2030
congressional seat allocations and their implications for the \textsc{bisect} redistricting
platform. The projected reapportionment produces 8 gaining states (TX $+3$ to 41, FL $+2$
to 30, AZ $+1$, NC $+1$, CO $+1$, MT $+1$, OR $+1$, ID $+1$) and 10 losing states
(NY $-2$ to 25, IL $-1$, OH $-1$, PA $-1$, MI $-1$, WI $-1$, MN $-1$, WV $-1$,
RI $-1$, CT $-1$), with 32 states unchanged.

The \textsc{bisect-apportion} crate handles projected populations directly, producing
seat allocations that match the Congressional Research Service's projected apportionment
within $\pm 1$ seat for 47 of 50 states. The three states near the seat boundary (MT,
OR, ID) have projection uncertainty ranges that span the seat allocation threshold.

We analyze algorithmic stability across the projected 2020$\to$2030 transition using
the county disruption metric from the T.9 paper: the fraction of county-seat assignments
that change between two redistricting cycles. Projected 2020$\to$2030 disruption averages
58\% across all states (lower than the 2010--2020 average of 66\%), driven by predictable
Sun Belt growth that makes optimal county groupings more stable. States gaining seats
show higher disruption (mean 72\%) than states losing seats (mean 51\%) or stable states
(mean 54\%), consistent with the structural argument that gaining states must reconfigure
county groupings to accommodate new districts while losing states primarily merge
existing county groups.

The Texas $k=41$ case is analyzed in detail. The prime-factor factorization of 41 is
trivial (41 is prime), so the \texttt{prime-factor} structure falls back to the binary
tree: a $20+21$ top-level split, with each half recursively bisected to 20 and 21 districts
respectively. We document the implications of this fallback for the ratio-optimal bisection
objective and compare to the \texttt{standard-bisect} and \texttt{nway} alternatives.
